\documentclass[12pt,a4paper]{article}
\usepackage[T1]{fontenc}
\usepackage[utf8]{inputenc}
\usepackage[spanish]{babel}
\usepackage{geometry}
\usepackage{booktabs}
\usepackage{xcolor}
\usepackage{listings}
\usepackage{hyperref}
\usepackage{enumitem}
\usepackage{titlesec}
\usepackage{parskip}

\geometry{margin=2.5cm}

\hypersetup{
  colorlinks=true,
  linkcolor=black,
  urlcolor=blue!70!black,
}

\lstdefinestyle{bash}{
  language=bash,
  basicstyle=\small\ttfamily,
  backgroundcolor=\color{gray!10},
  frame=single,
  rulecolor=\color{gray!40},
  breaklines=true,
  showstringspaces=false,
}

\title{%
  \textbf{Parlantes internos sin sonido}\\[0.4em]
  \large Diagnóstico y solución en el ASUS ROG Zephyrus G14 GA403UV\\
  con CachyOS + Omarchy (PipeWire / WirePlumber)
}
\author{pc-config}
\date{Septiembre de 2026}

\begin{document}

\maketitle
\tableofcontents
\newpage

% ─────────────────────────────────────────────
\section{Resumen}

Los parlantes internos del \textbf{ASUS ROG Zephyrus G14 GA403UV} dejaron de
sonar: PipeWire mostraba las barras de nivel moviéndose al reproducir audio
(por ejemplo, un video de YouTube), pero no salía sonido por los parlantes.

La causa no era el hardware, ni PipeWire, ni los amplificadores: el codec
\textbf{Realtek ALC285} quedó con los mixers \texttt{Speaker} en \textbf{0\%/off}
y \texttt{Bass Speaker} \textbf{off}, de modo que no salía señal hacia los
amplificadores \textbf{Cirrus Logic CS35L56}.

Peor aún: ese estado roto estaba guardado en
\texttt{/var/lib/alsa/asound.state}, por lo que
\texttt{alsa-restore.service} lo reaplicaba en \textbf{cada arranque}.

La solución es forzar los mixers, guardar el estado ALSA corregido y agregar un
servicio que los reaplique al arrancar y al volver de suspensión.

% ─────────────────────────────────────────────
\section{Entorno del sistema}

\begin{center}
\begin{tabular}{ll}
  \toprule
  Componente & Versión / valor \\
  \midrule
  Portátil            & ASUS ROG Zephyrus G14 GA403UV \\
  Sistema operativo   & CachyOS (base Arch Linux) \\
  Kernel activo       & \texttt{linux-cachyos} 7.0.9-1 \\
  Codec de audio      & Realtek ALC285 (tarjeta 2, HDA) \\
  Amplificadores      & 2 × Cirrus Logic CS35L56 (I2C) para altavoces \\
  Servidor de audio   & PipeWire 1.6.9 + WirePlumber 0.5.17 \\
  Utilidades ALSA     & \texttt{alsa-utils} \\
  \bottomrule
\end{tabular}
\end{center}

% ─────────────────────────────────────────────
\section{Síntoma}

\begin{itemize}[nosep]
  \item PipeWire/WirePlumber muestran el sink analógico como activo, sin mute y
        al 100\%, y las barras de nivel se mueven al reproducir.
  \item No se escucha nada por los parlantes internos (auriculares y HDMI no
        se vieron afectados).
  \item En \texttt{alsamixer} (o \texttt{amixer}) se observa:
        \texttt{Speaker} en 0\%/off y \texttt{Bass Speaker} off, mientras que
        \texttt{AMP1 Speaker} y \texttt{AMP2 Speaker} están en $\sim$89\%.
\end{itemize}

% ─────────────────────────────────────────────
\section{Diagnóstico}

El audio digital llega correctamente hasta los amplificadores: durante la
reproducción, los dispositivos I2C de los CS35L56 pasan a estado
\texttt{active} y los controles \texttt{AMP1/AMP2 Speaker} están en su nivel
calibrado (400/448, 0.00\,dB).

El problema está en la etapa anterior: el ALC285 tiene dos controles que
funcionan como compuerta de los parlantes:

\begin{lstlisting}[style=bash]
$ amixer -c 2 sget Speaker
  Front Left: Playback 0 [0%] [-65.25dB] [off]
  Front Right: Playback 0 [0%] [-65.25dB] [off]

$ amixer -c 2 sget 'Bass Speaker'
  Front Left: Playback [off]
  Front Right: Playback [off]
\end{lstlisting}

Además, el estado guardado de ALSA contenía esos valores rotos:

\begin{lstlisting}[style=bash]
awk '/^state.Generic_1/,/^state.Controller/' /var/lib/alsa/asound.state \
  | grep -A2 -E "name '(Speaker|Bass Speaker)"
# value.0 0 / value.1 0          <- Speaker Playback Volume
# value.0 false / value.1 false  <- Speaker Playback Switch
# value.0 false / value.1 false  <- Bass Speaker Playback Switch
\end{lstlisting}

Por eso \texttt{alsa-restore.service} volvía a dejar los parlantes mudos en
cada arranque.

% ─────────────────────────────────────────────
\section{Solución}

\subsection*{Aplicar y guardar}

\begin{lstlisting}[style=bash]
amixer -c 2 sset Speaker 100% unmute
amixer -c 2 sset 'Bass Speaker' on
sudo alsactl store
\end{lstlisting}

\subsection*{Persistencia con systemd}

El paso 19 (\texttt{setup-speaker-fix.sh}) instala:

\begin{itemize}[nosep]
  \item \texttt{/usr/local/bin/asus-g14-speaker-fix} — detecta la tarjeta del
        ALC285 y fuerza ambos mixers. No usa el número de tarjeta fijo: lo
        obtiene de \texttt{aplay -l} (con fallback por control
        \texttt{Bass Speaker}).
  \item \texttt{/etc/systemd/system/asus-g14-speaker-fix.service} —
        \texttt{Type=oneshot}, \texttt{After=alsa-restore.service} y
        \texttt{WantedBy=multi-user.target suspend.target hibernate.target
        hybrid-sleep.target suspend-then-hibernate.target}, para que corra al
        arrancar y al volver de suspensión aunque el estado guardado se haya
        vuelto a romper.
\end{itemize}

Verificación:

\begin{lstlisting}[style=bash]
systemctl status asus-g14-speaker-fix.service
amixer -c 2 sget Speaker         # 100% y on
amixer -c 2 sget 'Bass Speaker'  # on
\end{lstlisting}

% ─────────────────────────────────────────────
\section{Notas}

\begin{itemize}[nosep]
  \item No tocar \texttt{AMP1 Speaker} / \texttt{AMP2 Speaker} por encima de
        0.00\,dB: forman parte de la protección de los parlantes.
  \item Si en el futuro se cambia el perfil de la tarjeta o se conectan
        auriculares y el problema reaparece, ejecutar
        \texttt{asus-g14-speaker-fix} (o reiniciar el servicio).
  \item El micrófono interno es un problema distinto: ver
        \texttt{microfono-interno.pdf}.
\end{itemize}

\end{document}
